% Generated draft: requires mathematical and source review.
\documentclass[11pt]{article}
\usepackage[margin=25mm]{geometry}
\usepackage{fontspec}
\setmainfont{Latin Modern Roman}
\usepackage{amsmath,amssymb,mathtools,unicode-math}
\setmathfont{Latin Modern Math}
\usepackage{enumitem,needspace,xurl,hyperref}
\setlength{\parindent}{0pt}
\setlength{\parskip}{5pt}
\setlength{\emergencystretch}{3em}
\newcommand{\N}{\mathbb{N}}
\newcommand{\Z}{\mathbb{Z}}
\newcommand{\Q}{\mathbb{Q}}
\newcommand{\R}{\mathbb{R}}
\newcommand{\C}{\mathbb{C}}
\newcommand{\F}{\mathbb{F}}
\newcommand{\A}{\mathbb{A}}
\newcommand{\PP}{\mathbb{P}}
\newcommand{\E}{\mathbb{E}}
\newcommand{\eps}{\varepsilon}
\newcommand{\dd}{\,\mathrm{d}}
\begin{document}
\section*{UnsolvedMath ID: 30005460}
\textbf{Source problem code:} catalogue entry
\subsection{Definitions and mathematical statement}
For a real topological vector space \(V\), a subset \(C\subseteq V\) is a closed convex cone if it is nonempty, closed in the specified topology of \(V\), and satisfies
\[
\forall f,g\in C\;\forall a,b\in\R_{\geq 0},\qquad af+bg\in C.
\]
This is a general definition only; the supplied record does not identify the ambient space or topology for the set below.

The source question is:
\begin{quote}
For the source-intended admissible parameters \(n,m,k\), is \(\Sigma_{n,m}(2k+1)\) a closed convex cone?
\end{quote}
The supplied curation describes \(\Sigma_{n,m}(2k+1)\) as an odd-power sum-of-squares-root set, but gives no membership definition. In particular, it does not specify the parameter ranges, whether the underlying objects are polynomials or homogeneous forms, the degree convention, or the ambient topology. Consequently, the notation cannot be expanded into a fully specified mathematical assertion from the supplied evidence alone.

The supplied literature assessment attributes closedness to the source and describes convexity as difficult and conjectural. Thus closedness is presented in the curation as established, while convexity remains the stated unresolved issue. These are supplied source assessments, not independently verified mathematical conclusions.
\subsection{Short English statement}
The question asks whether the source's odd-power sum-of-squares-root set \(\Sigma_{n,m}(2k+1)\) is a closed convex cone. The supplied assessment says closedness is proved and convexity remains open. A precise formulation requires the missing definition of this set, its admissible parameters, and its ambient space and topology.
\subsection{Sources}
\begin{itemize}
\item Supplied local UnsolvedMath registry record, permanent ID \(30005460\), source problem code \texttt{catalogue entry}, title ``Convexity of Odd-Power Sum-of-Squares Cones.'' The registry labels the problem \texttt{open} and the statement \texttt{corrected\_verified}; these are supplied metadata labels. Its verification note reports direct checking of the official report and restoration of definitions, but those definitions are absent from the supplied record.
\item \emph{Real Algebraic Geometry with a View toward Koopman Operator Methods} (2023), Oberwolfach Reports, pp.\ 741--811, DOI \texttt{10.4171/owr/2023/14}. The extracted question carries page number 779. Supplied source URL and literature-source link, labeled ``original Oberwolfach report'': \url{https://doi.org/10.4171/owr/2023/14}. No individual proposer is identified in the supplied record.
\item Literature assessment dated exactly \texttt{2026-08-22}: the supplied curation says the source proves closedness, leaves convexity difficult and conjectural, and that no primary result found in that review establishes or refutes convexity for the stated general set. This is supplied dated curation, not independently verified. The background explicitly requests independent verification of time-sensitive status and cited claims.
\end{itemize}
\end{document}
